\section{Conclusion}
\label{sec:conclusion}

The DDM fits the class's response-time curve more closely than my own.
My choices are relatively sensitive to bid differences, yet the bins where
I take longest are not those with the closest bids. This matters because
the time curve uses the sensitivity estimated from choices; changing $A$
and $t_0$ cannot move its peak away from equal values. The class means do
peak near equal bids.

Both fits depend on how well the bids measure value at the moment of choice.
Bids may contain noise or reflect values that changed before the choice
task. In the pooled data, averaging across students with different choice
sensitivities could also help explain a steeper center and less extreme
tails than a single logistic curve predicts. The close fit to class
response times does not establish that everyone shares the same
$\alpha$, $A$, and $t_0$.
